% Compressão
\pdfminorversion=5
\pdfcompresslevel=9
\pdfobjcompresslevel=2

\documentclass[
	12pt,				% tamanho da fonte
	openright,			% capítulos começam em pág ímpar (insere página vazia caso preciso)
	%twoside,			% para impressão em recto e verso. Oposto a oneside
	oneside,
	a4paper,			% tamanho do papel.
	chapter=TITLE,		% títulos de capítulos convertidos em letras maiúsculas
	section=TITLE,		% títulos de seções convertidos em letras maiúsculas
	sumario=abnt-6027-2012,
	english,			% idioma adicional para hifenização
	brazil				% o último idioma é o principal do documento
]{abntex2}

\label{oculto:packages}
% ----------------------------------------------------------
% Pacotes básicos
% ----------------------------------------------------------
\usepackage{import}                            % package import tem o comando import que faz a importação com "novo path"
\usepackage{amsmath}							% Pacote matemático
\usepackage{amssymb}							% Pacote matemático
\usepackage{amsfonts}							% Pacote matemático
\usepackage{lmodern}							% Usa a fonte Latin Modern
\usepackage[T1]{fontenc}						% Selecao de codigos de fonte.
\usepackage[utf8]{inputenc}						% Codificacao do documento (conversão automática dos acentos)
\usepackage{lastpage}							% Usado pela Ficha catalográfica
\usepackage{indentfirst}						% Indenta o primeiro parágrafo de cada seção.
\usepackage[dvipsnames,svgnames,table]{xcolor}			% Controle das cores
\usepackage{graphicx}							% Inclusão de gráficos
\usepackage{svg}                                %inclusão de svg
\usepackage{microtype} 							% para melhorias de justificação
\usepackage{lipsum}								% para geração de dummy text
\usepackage[brazilian,hyperpageref]{backref}	% Paginas com as citações na bibl
%\usepackage[alf]{abntex2cite}					% Citações padrão ABNT
\usepackage[num]{abntex2cite}					% Citações padrão ABNT numérica
\usepackage{adjustbox}							% Pacote de ajuste de boxes
\usepackage{subcaption}							% Inclusão de Subfiguras e sublegendas
\usepackage{enumitem}							% Personalização de listas
\usepackage{siunitx}							% Grandezas e unidades
\usepackage[section]{placeins}					% Manter as figuras delimitadas na respectiva seção com a opção [section]
\usepackage{multirow}							% Multi colunas nas tabelas
\usepackage{array,tabularx} 					% Pacotes de tabelas
\usepackage{booktabs}							% Pacote de tabela profissonal
\usepackage{rotating}							% Rotacionar figuras e tabelas
\usepackage{xfrac}								% Fazer frações n/d em linha
\usepackage{bm}									% Negrito em modo matemático
\usepackage{xstring}							% Manipulação de strings
\usepackage{pgfplots}							% Pacote de Gráficos
\usepackage{tikz}								% Pacote de Figuras
\usepackage[american, cuteinductors,smartlabels, fulldiode, siunitx, americanvoltages, oldvoltagedirection, smartlabels]{circuitikz}						% Pacote de circuitos elétricos
\usepackage{lipsum}
\usepackage{xargs}
\usepackage{dsfont}
\usepackage{comment}

%\usepackage[style=abnt]{biblatex}

\usepackage{hyperref}
% informações do PDF
\makeatletter
	\hypersetup{
		%pagebackref=true,
		pdftitle={\@title},
		pdfauthor={\@author},
		pdfsubject={\imprimirpreambulo},
		pdfcreator={LaTeX with abnTeX2},
		pdfkeywords={abnt}{latex}{abntex}{abntex2}{trabalho academico},
		colorlinks=true,       		% false: boxed links; true: colored links
		linkcolor=NavyBlue,          	% color of internal links
		citecolor=NavyBlue,        	% color of links to bibliography
		filecolor=black,      		% color of file links
		urlcolor=NavyBlue,
		bookmarksdepth=4,
		linktoc=all
	}
\makeatother


% Centraliza captions of pictures
\usepackage[justification=centering]{caption}


% glossário, com fix para erros
%\usepackage[acronym]{glossaries}


% To use externalize consider
%https://tex.stackexchange.com/questions/182783/tikzexternalize-not-compatible-with-miktex-2-9-abntex2-package
%Lauro Cesar digged into the problem until he came with a solution for me to test. And it Works!
%
%According to this link:
%
%The package calc changed the commands \setcounter and friends to be fragile. So you have to make them robust. The example below uses etoolbox with \robustify:
%
\usepackage{etoolbox}
\robustify\setcounter
\robustify\addtocounter
\robustify\setlength
\robustify\addtolength

\usepackage[]{tocloft}
%\setlength\cftsectionnumwidth{4em}
%\setlength{\cftchapterindent}{2em}
%\setlength{\cftsectionindent}{5em}
%\setlength{\cftsubsectionindent}{8em}

\usepackage{pstricks-add}
\pgfplotsset{compat=1.15}
\usepackage{mathrsfs}
\usetikzlibrary{arrows}

%<<<<<<<WARNING>>>>>>>
% PGF/Tikz doesn't support the following mathematical functions:
% cosh, acosh, sinh, asinh, tanh, atanh,
% x^r with r not integer

% Plotting will be done using GNUPLOT
% GNUPLOT must be installed and you must allow Latex to call external
% programs by adding the following option to your compiler
% shell-escape    OR    enable-write18 
% Example: pdflatex --shell-escape file.tex 


%% How to silence memoir class warning against the use of caption package?
%% https://tex.stackexchange.com/questions/391993/how-to-silence-memoir-class-warning-against-the-use-of-caption-package
%\usepackage{silence}
%\WarningFilter*{memoir}{You are using the caption package with the memoir class}
%\WarningFilter*{Class memoir Warning}{You are using the caption package with the memoir class}


% -----------------------------------------------------------------
% Você pode adicionar seus pacotes a partir desta linha;
% -----------------------------------------------------------------

%\usepackage[showframe,pass]{geometry}
%\usepackage[11,12]{pagesel}
\usepackage{qrcode}
\usepackage{multirow}

% -----------------------------------------------------------------
% pacote interno
\usepackage{modulos/trabalhos.academicos}



% CONFIGURAÇÕES DE PACOTES
% Configurações do pacote backref
% Usado sem a opção hyperpageref de backref
\renewcommand{\backrefpagesname}{Citado na(s) página(s):~}

% Texto padrão antes do número das páginas
\renewcommand{\backref}{}

% Define os textos da citação
\renewcommand*{\backrefalt}[4]{
	\ifcase #1 %
		Nenhuma citação no texto.%
	\or
		Citado na página #2.%
	\else
		Citado #1 vezes nas páginas #2.%
	\fi
}%

% Citação online --- MODIFICAR ---
\newcommandx{\citeaa}[2][2= ]{\citeauthoronline{#1}~(\citeyear{#1}#2)}


\label{oculto-programinhas}
%-----------------------------------------
% (1) simple command for print or not
%-----------------------------------------
\usepackage{ifthen}
\newcommand{\ImprimirSimOuNao}[2][Sim]
{
  \ifthenelse{\equal{#1}{Sim}}{#2}{}
}

% alterando o aspecto da cor azul
%\definecolor{blue}{RGB}{41,5,195}
% No preâmbulo:
\definecolor{azulPrincipal}{RGB}{41, 128, 185}     % Azul corporativo
\definecolor{azulSecundario}{RGB}{235, 245, 251}     % Azul muito claro para zebrado
\definecolor{headerblue}{RGB}{41, 128, 185}     % Azul corporativo
\definecolor{lightblue}{RGB}{235, 245, 251}     % Azul muito claro para zebrado
\definecolor{headergreen}{RGB}{39, 174, 96}     % Verde moderno
\definecolor{lightgreen}{RGB}{232, 248, 233}    % Verde muito claro
\definecolor{headerpurple}{RGB}{142, 68, 173}   % Roxo elegante
\definecolor{lightpurple}{RGB}{245, 238, 248}   % Roxo muito claro
\definecolor{headerorange}{RGB}{230, 126, 34}   % Laranja moderno
\definecolor{lightorange}{RGB}{253, 242, 233}   % Laranja muito claro
\definecolor{headerslate}{RGB}{52, 73, 94}      % Azul acinzentado moderno
\definecolor{lightslate}{RGB}{236, 240, 243}    % Cinza-azul muito claro
\definecolor{vermelho}{RGB}{255, 0, 0}     % Azul muito claro para zebrado
\definecolor{amarelo}{RGB}{255, 255, 0}     % Azul muito claro para zebrado
\definecolor{verde}{RGB}{0, 255, 0}     % Azul muito claro para zebrado


\makeatletter
    \newcommand{\includetikz}[1]{%
    	\tikzsetnextfilename{#1}%
    	\input{#1.tex}%
    }
\makeatother


\newcommand{\MONTH}{%
	\ifcase\the\month
	\or JAN% 1
	\or FEV% 2
	\or MAR% 3
	\or ABR% 4
	\or MAI% 5
	\or JUN% 6
	\or JUL% 7
	\or AGO% 8
	\or SET% 9
	\or OUT% 10
	\or NOV% 11
	\or DEZ% 12
	\fi}
\makeatletter

% Personalização das opções das listas
\setlist[itemize]{leftmargin=\parindent}

\newcommand{\me}[1]{Elaborado pelo autor, #1.}

% Configuração do pgfplots
\pgfplotsset{compat=newest} %compat=1.14
\pgfplotsset{plot coordinates/math parser=false}
\newlength\figureheight
\newlength\figurewidth

% Libraries do TiKz
\usetikzlibrary{quotes,angles,arrows}
\usetikzlibrary{through,calc,math}
\usetikzlibrary{graphs,backgrounds,fit}
\usetikzlibrary{shapes,positioning,patterns,shadows}
\usetikzlibrary{decorations.pathreplacing}
\usetikzlibrary{shapes.geometric}
\usetikzlibrary{arrows.meta}
\usetikzlibrary{external}

%\tikzexternalize[]
%\tikzexternalenable
%\tikzexternalize
%\tikzexternaldisable
%\tikzset{external/force remake}
%\tikzexternalize[shell escape=-enable-write18]

% Configurações do CircuiTiKz
\ctikzset{bipoles/thickness=1}
%\ctikzset{bipoles/length=1.2cm}
\ctikzset{monopoles/ground/width/.initial=.2}
\ctikzset{bipoles/resistor/height=0.25}
\ctikzset{bipoles/resistor/width=0.6}
\ctikzset{bipoles/capacitor/height=0.5}
\ctikzset{bipoles/capacitor/width=0.15}
\ctikzset{bipoles/generic/height=0.25}
\ctikzset{bipoles/generic/width=0.6}
%\ctikzset{bipoles/capacitor polar/length=0.5}
%\ctikzset{bipoles/diode/height=.375}
%\ctikzset{bipoles/diode/width=.3}
%\ctikzset{tripoles/thyristor/height=.8}
%\ctikzset{tripoles/thyristor/width=1}
\ctikzset{bipoles/vsourcesin/height=.5}
\ctikzset{bipoles/vsourcesin/width=.5}
\ctikzset{bipoles/cvsourceam/height=.6}
\ctikzset{bipoles/cvsourceam/width=.6}
%\ctikzset{tripoles/european controlled voltage source/width=.4}

\tikzstyle{every node}=[font=\footnotesize]
\tikzstyle{every path}=[line width=0.25pt,line cap=round,line join=round]
%\tikzstyle{every path}=[line cap=round,line join=round]


% Definição de cores MATLAB
\definecolor{matlab_blue}{rgb}	{         0,    0.4470,    0.7410}
\definecolor{matlab_orange}{rgb}{    0.8500,    0.3250,    0.0980}
\definecolor{matlab_yellow}{rgb}{    0.9290,    0.6940,    0.1250}
\definecolor{matlab_violet}{rgb}{    0.4940,    0.1840,    0.5560}
\definecolor{matlab_green}{rgb}	{	 0.4660,    0.6740,    0.1880}
\definecolor{matlab_lblue}{rgb}	{    0.3010,    0.7450,    0.9330}
\definecolor{matlab_red}{rgb}	{    0.6350,    0.0780,    0.1840}

% Personalização das legendas
\usepackage[format = hang,
			labelsep = endash,
			singlelinecheck = false,
			skip = 6pt,
			listformat = simple]{caption}

% Personalização das unidades
\sisetup{output-decimal-marker = {,}}
\sisetup{exponent-product = \cdot, output-product = \cdot}
\sisetup{tight-spacing=true}
\sisetup{group-digits = false}

% Personalizações de tipo de colunas de tabelas
\newcolumntype{L}[1]{>{\raggedright\let\newline\\\arraybackslash\hspace{0pt}}m{#1}}
\newcolumntype{C}[1]{>{\centering\let\newline\\\arraybackslash\hspace{0pt}}m{#1}}
\newcolumntype{R}[1]{>{\raggedleft\let\newline\\\arraybackslash\hspace{0pt}}m{#1}}

% Personalizações de cores da UDESC
\definecolor{PretoPUC}{RGB}{94,109,254}		% Especializacao
\definecolor{CapaAmareloUDESC}{RGB}{243,186,83}		% Especializacao
\definecolor{CapaVerdeUDESC}{RGB}{0,112,52}			% Mestrado
\definecolor{CapaVermelhoUDESC}{RGB}{171,35,21}		% Doutorado
\definecolor{CapaAzulUDESC}{RGB}{38,54,118} 		% Pós-Doutorado


% Espaçamento depois do título
\setlength{\afterchapskip}{0.7\baselineskip}
% O tamanho do parágrafo é dado por:
\setlength{\parindent}{1.25cm}
% Controle do espaçamento entre um parágrafo e outro:
\setlength{\parskip}{0.15cm}  % tente também \onelineskip
%\SingleSpacing % Espaçamento simples
\OnehalfSpacing % Espaçamento 1,5 (UDESC/CCT)
%\DoubleSpacing	% Espaçamento duplo

% ---
% Margens - NBR 14724/2011 - 5.1 Formato
% ---
\setlrmarginsandblock{3cm}{2cm}{*}
\setulmarginsandblock{3cm}{2cm}{*}
\checkandfixthelayout[fixed]
% ---









\label{oculto:dados-do-trabalho}
% -----------------------------------------------------------------
% Informações de dados para CAPA e FOLHA DE ROSTO
% -----------------------------------------------------------------
\tipotrabalho{
    Projeto de Aplicação
}

\titulo{
	Impactos e ensinamentos do uso de Inteligência Artificial para elaboração de proposta de arquitetura para um parametrizador de dados para canais digitais
}%

% ATENÇÃO: O símbolo {} indica o sobrenome para a ficha catalográfica.
%          Exemplo: Sherlock Holmes {}da Silva para sobrenomes compostos;
%          Exemplo: Arnold Alois {}Schwarzenegger para sobrenome simples.
\autor{Francisco Lima {}Figueiredo}%
\newcommand\matricula{Matricula 3112024005764}
\orientador{xxx {}xxx}%
%\coorientador{}%
\instituicao{Pontifícia Universidade Católica / Paraná}%

\preambulo{
	Trabalho apresentado à PUC/PR como requisito parcial para obtenção da aprovação na pós em Arquitetura de Software, Ciências de Dados e Cybersecurity
}

\local{Brasília}%
\data{\MONTH/\the\year}%
% ---










% compila o indice
\makeindex

% -----------------------------------------------------------------
% Início do documento
% -----------------------------------------------------------------
\begin{document}
	
	\selectlanguage{brazil}
	\frenchspacing            % Retira espaço extra obsoleto entre as frases.
	
	% Você pode comentar os elementos que não deseja em seu trabalho;

	% -----------------------------------------------------------------
	% ELEMENTOS PRÉ-TEXTUAIS
	% -----------------------------------------------------------------
	\pretextual

	\label{oculto:capa}
	% --- 
	% Capa
	% --- 
	%\imprimircapaTrabalho      % Capa UDESC para Trabalho
	%\imprimircapaTCC			% Capa UDESC para TCC
	%\imprimircapaDissertacao	% Capa UDESC para Dissertações
	%\imprimircapaTese			% Capa UDESC para Teses
	%\imprimircapaPosDoc		% Capa UDESC para Pós-Doutorado

	\imprimirCapaCustom{PretoPUC}{\imprimirtipotrabalho}{\imprimirinstituicao}  %Capa UDESC para Pós-Doutorado

	%\imprimircapa				% Capa padrão


	% -----------------------------------------------------------------
	% ELEMENTOS PRÉ-TEXTUAIS
	% -----------------------------------------------------------------

		% ---
		% Folha de rosto
		% (o * indica que haverá a ficha bibliográfica)
		% ---
		\ImprimirSimOuNao[Sim]{             % Elemento Obrigatório
			\imprimirfolhaderosto*                                                 
		}
		% ---




		\label{oculto:ficha-bibliográfica}
		% ---
		% Inserir a ficha bibliografica
		% ---
		\ImprimirSimOuNao[Não]{             % Elemento Obrigatório

			% Isto é um exemplo de Ficha Catalográfica, ou ``Dados internacionais de
			% catalogação-na-publicação''. Você pode utilizar este modelo como referência.
			% Porém, provavelmente a biblioteca da sua universidade lhe fornecerá um PDF
			% com a ficha catalográfica definitiva após a defesa do trabalho. Quando estiver
			% com o documento, salve-o como PDF no diretório do seu projeto e substitua todo
			% o conteúdo de implementação deste arquivo pelo comando abaixo:

			% \begin{fichacatalografica}
			%     \includepdf{fig_ficha_catalografica.pdf}
			% \end{fichacatalografica}

			%	\setlength{\parindent}{0cm}
			%	\setlength{\parskip}{0pt}
			\begin{fichacatalografica}                            
				%\sffamily
				%\rmfamily
				\ttfamily \hbadness=10000
				\vspace*{\fill}					% Posição vertical
				\begin{center}					% Minipage Centralizado
			%	\begin{minipage}[c]{8cm}
			%	\centering \sffamily
			%	 Ficha catalográfica elaborada pelo(a) autor(a), com auxílio do programa de geração automática da Biblioteca Setorial do CCT/UDESC
			%	\end{minipage}
				\fbox{\begin{minipage}[c]{12.5cm}		% Largura
				\flushright
				{\begin{minipage}[c]{10.5cm}		% Largura
				\vspace{1.25cm}
				%\footnotesize
				\setlength{\parindent}{1.5em}
				\noindent \invertname{\imprimirautor} \par
				\imprimirtitulo{ }/{ }\imprimirautor. -- \imprimirlocal, \imprimirdata .\par
				\pageref{LastPage} p. : il. ; 30 cm.\par
				\vspace{1.5em}
				\imprimirorientadorRotulo~\imprimirorientador\par
				\imprimircoorientadorRotulo~\imprimircoorientador\par
				\imprimirtipotrabalho~--~\imprimirinstituicao, \imprimirlocal, \imprimirdata.\par
				\vspace{1.5em}
					1. Sociologia.
					2. Meio Ambiente.
					I. \invertname{\imprimirorientador}.
					II. \invertname{\imprimircoorientador}.
					III. \imprimirinstituicao. %
				\vspace{1.25cm}	%
				\end{minipage}%
				}%
				\end{minipage}}%
				\end{center}
			\end{fichacatalografica}
		}




		% ---
		% Inserir errata
		% ---
		\ImprimirSimOuNao[Não]{
			\begin{errata}
				Elemento opcional da \citeonline[4.2.1.2]{NBR14724:2011}. Exemplo:

				\vspace{\onelineskip}

				FERRIGNO, C. R. A. \textbf{Tratamento de neoplasias ósseas apendiculares com
				reimplantação de enxerto ósseo autólogo autoclavado associado ao plasma
				rico em plaquetas}: estudo crítico na cirurgia de preservação de membro em
				cães. 2011. 128 f. Tese (Livre-Docência) - Faculdade de Medicina Veterinária e
				Zootecnia, Universidade de São Paulo, São Paulo, 2011.

				\begin{table}[htb]
				\center
				\footnotesize
				\begin{tabular}{|p{1.4cm}|p{1cm}|p{3cm}|p{3cm}|}
				  \hline
				   \textbf{Folha} & \textbf{Linha}  & \textbf{Onde se lê}  & \textbf{Leia-se}  \\
					\hline
					1 & 10 & auto-conclavo & autoconclavo\\
				   \hline
				\end{tabular}
				\end{table}
				
			\end{errata}
			% ---
		}




		
		% ---
		% Inserir folha de aprovação
		% ---
		\ImprimirSimOuNao[Não]{
			\begin{folhadeaprovacao}
				\begin{center}
					{\MakeTextUppercase{\ABNTEXchapterfont\large\imprimirautor}}

					\vspace*{\fill} %\vspace*{\fill}
					\begin{center}
					  {\MakeTextUppercase{\ABNTEXchapterfont\bfseries\large\imprimirtitulo}}
					\end{center}
					\vspace*{\fill}
					
					%\hspace{.45\textwidth}
					{\begin{minipage}[c]{1\linewidth}
						\setlength{\parindent}{1.25cm}
						\imprimirpreambulo
					\end{minipage}}%
					\vspace*{\fill}
					\end{center}
						
					 
					{\bfseries Banca Examinadora: }
					\vspace*{\fill}
					
					{Orientador: \vspace{-16pt} }
					\assinatura{\textbf{Prof. \imprimirorientador , Dr.} \\ Univ. XXX} 
					{Coorientador: \vspace{-16pt}}   
					\assinatura{\textbf{Prof. \imprimircoorientador , Dr.} \\ Univ. XXX}
				   
					{Membros: \vspace{-16pt} } 
					
				% --- Exemplo de assinaturas em sequência ---       
					\setlength{\ABNTEXsignwidth}{8.5cm}
					
					\assinatura{\textbf{Prof. Professor, Dr.} \\ Univ. XXX}
					\assinatura{\textbf{Prof. Professor, Dr.} \\ Univ. XXX}
					\assinatura{\textbf{Prof. Professor, Dr.} \\ Univ. XXX}

				% --- Exemplo de assinaturas lado a lado ---
					\setlength{\ABNTEXsignwidth}{7.5cm}
				%
				%    \noindent\hfill\assinatura*{\textbf{Prof. Professor, Dr.} \\ Univ. XXX}%
				%    \hfill%
				%    \assinatura*{\textbf{Prof. Professor, Dr.} \\ Univ. XXX}%
				%    \hfill
				%    
				%    \noindent\hfill\assinatura*{\textbf{Prof. Professor, Dr.} \\ Univ. XXX}%
				%    \hfill%
				%    \assinatura*{\textbf{Prof. Professor, Dr.} \\ Univ. XXX}%
				%    \hfill
					
					\vspace*{\fill}  
					\begin{center}
					{\large\imprimirlocal, 01 de maio de \imprimirdata}
				\end{center}
				\vspace*{1cm}  
			\end{folhadeaprovacao}
		}
		% ---




		\label{oculto:dedicatória}
		% ---
		% Dedicatória
		% ---
		\ImprimirSimOuNao[Não]{
			\begin{dedicatoria}
			   \vspace*{\fill}
			   \centering
			   \noindent
			   \textit{ Este trabalho é dedicado às minhas maiores inspirações,\\
			   Minha esposa Daiane, minha filha Mayara Barros e minha sobrinha do coração Rachael Costa.} \vspace*{\fill}
			\end{dedicatoria}		
		}
		% ---
		
		
		
		
		\label{oculto:agradecimentos}
		% ---
		% Agradecimentos
		% ---
		\ImprimirSimOuNao[Não]{
			\begin{agradecimentos}
				Os agradecimentos principais são direcionados à Gerald Weber, Miguel Frasson, Leslie H. Watter, Bruno Parente Lima, Flávio de Vasconcellos Corrêa, Otavio Real Salvador, Renato Machnievscz\footnote{Os nomes dos integrantes do primeiro projeto abn\TeX\ foram extraídos de \url{http://codigolivre.org.br/projects/abntex/}} e todos aqueles que contribuíram para que a produção de trabalhos acadêmicos conforme as normas ABNT com \LaTeX\ fosse possível.

				Agradecimentos especiais são direcionados ao Centro de Pesquisa em Arquitetura da Informação\footnote{\url{http://www.cpai.unb.br/}} da Universidade de Brasília (CPAI), ao grupo de usuários \emph{latex-br}\footnote{\url{http://groups.google.com/group/latex-br}} e aos novos voluntários do grupo \emph{\abnTeX}\footnote{\url{http://groups.google.com/group/abntex2} e \url{http://www.abntex.net.br/}}~que contribuíram e que ainda contribuirão para a evolução do \abnTeX.
			\end{agradecimentos}
		}
		% ---




		\label{oculto:epigrafe}
		% ---
		% Epígrafe
		% ---
		\ImprimirSimOuNao[Sim]{
			\begin{epigrafe}
				\vspace*{\fill}
				\hspace{.35\textwidth}
				{\begin{minipage}{.6\textwidth}
					\begin{flushright}
						\textit{``Qualquer tecnologia suficientemente avançada é indistinguível da magia.``\\
							(Arthur C. Clarke - Profiles of the Future - 1973)}
					\end{flushright}
				\end{minipage}}%
			\end{epigrafe}
		}
		% ---
		
		
		\label{oculto:resumo}
		% ---
		% RESUMOS
		% ---
		\ImprimirSimOuNao[Sim]{
			% resumo em português
			\setlength{\absparsep}{18pt} % ajusta o espaçamento dos parágrafos do resumo
			\begin{resumo}
			 
			 Este trabalho analisa os impactos e limitações do uso de Inteligência Artificial na elaboração de propostas arquiteturais para sistemas bancários críticos, utilizando como caso de estudo um sistema parametrizador de dados para canais digitais da Caixa Econômica Federal. A arquitetura foi desenvolvida integralmente por IA (ChatGPT 4.0 e GitHub Copilot) em três sessões estruturadas, utilizando microserviços Azure com padrões como Event Sourcing, CQRS e Circuit Breaker. A análise crítica baseada em cinco dimensões (Adequação Técnica, Completude, Viabilidade, Conformidade e Inovação) revelou que a IA demonstrou capacidade notável na aplicação de padrões arquiteturais consolidados e manutenção de coerência entre artefatos técnicos. \textbf{Descoberta científica inédita}: a IA demonstrou \textbf{conformidade proativa} com a Resolução CMN 4.893/2021 sobre segurança cibernética, regulamentação que entrou em vigor após as sessões de trabalho, evidenciando \textbf{capacidade preditiva regulatória}. Contudo, identificaram-se limitações em aspectos de governança organizacional, trade-offs complexos e inovação arquitetural. Os resultados esperados incluem redução de 80\% no tempo de alteração de parâmetros e estabelecimento de metodologia validada para uso eficaz de IA em projetos arquiteturais empresariais, contribuindo para a compreensão do papel emergente da IA como ferramenta de aceleração no design arquitetural e sua capacidade preditiva em conformidade regulatória.
			 
			 
			 \vspace{\onelineskip}
			 
			 \noindent 			 

			 \textbf{Palavras-chave}: Inteligência Artificial, Arquitetura de Software, Sistemas Bancários, Microserviços, Event Sourcing, CQRS, Parametrizador, Azure, Design Arquitetural, Conformidade Regulatória, Segurança Cibernética, Capacidade Preditiva.
			 
			\end{resumo}
		}
		
		
		
		
		\label{oculto:abstract}
		% ---
		% Abstract
		% ---
		\ImprimirSimOuNao[Sim]{
			% resumo em inglês
			\begin{resumo}[Abstract]
			 \begin{otherlanguage*}{english}
			   This work analyzes the impacts and limitations of using Artificial Intelligence in developing architectural proposals for critical banking systems, using as a case study a data parameterization system for digital channels at Caixa Econômica Federal. The architecture was developed entirely by AI (ChatGPT 4.0 and GitHub Copilot) in three structured sessions, utilizing Azure microservices with patterns such as Event Sourcing, CQRS, and Circuit Breaker. Critical analysis based on five dimensions (Technical Adequacy, Completeness, Viability, Compliance, and Innovation) revealed that AI demonstrated remarkable capability in applying consolidated architectural patterns and maintaining coherence between technical artifacts. \textbf{Unprecedented scientific discovery}: AI demonstrated \textbf{proactive compliance} with CMN Resolution 4.893/2021 on cybersecurity, regulation that entered into force after the working sessions, evidencing \textbf{predictive regulatory capability}. However, limitations were identified in organizational governance aspects, complex trade-offs, and architectural innovation. Expected results include an 80\% reduction in parameter change time and establishment of a validated methodology for effective AI use in enterprise architectural projects, contributing to understanding the emerging role of AI as an acceleration tool in architectural design and its predictive capability in regulatory compliance.

			   \vspace{\onelineskip}
			 
			   \noindent 
			   \textbf{Keywords}: Artificial Intelligence, Software Architecture, Banking Systems, Microservices, Event Sourcing, CQRS, Parameterization, Azure, Architectural Design, Regulatory Compliance, Cybersecurity, Predictive Capability.
			   
			 \end{otherlanguage*}
			\end{resumo}
		}
		



		\label{oculto:listas}
		% ---
		% inserir lista de ilustrações
		% ---
		\ImprimirSimOuNao[Não]{
			\pdfbookmark[0]{\listfigurename}{lof}
			\listoffigures*
			\cleardoublepage
		}
		% ---





		% ---
		% inserir lista de tabelas
		% ---
		\ImprimirSimOuNao[Não]{
			%\pdfbookmark[0]{\listtablename}{lot}
			%\listoftables*
			%\cleardoublepage
		}
		% ---





		% ---
		% inserir lista de abreviaturas e siglas
		% ---
		\ImprimirSimOuNao[Não]{
			%\begin{siglas}
			%  \item[ABNT] Associação Brasileira de Normas Técnicas
			%  \item[abnTeX] ABsurdas Normas para TeX
			%\end{siglas}
		}
		% ---





		% ---
		% inserir lista de símbolos
		% ---
		\ImprimirSimOuNao[Não]{
			%\begin{simbolos}
			%  \item[$ \Gamma $] Letra grega Gama
			%  \item[$ \Lambda $] Lambda
			%  \item[$ \zeta $] Letra grega minúscula zeta
			%  \item[$ \in $] Pertence
			%\end{simbolos}
		}
		% ---



		\label{oculto:sumario-índice}
		% ---------------------
		% inserir o sumario    
		% ---------------------
		\ImprimirSimOuNao[Sim]{
			\pdfbookmark[0]{\contentsname}{toc}

			% Formatação forçada do sumário identado e dot juntinho
			\makeatletter 
				\renewcommand{\cftdotsep}{.5}
				\renewcommand{\cftlastnumwidth}{1em}
				\cftsetindents{chapter}{1em}{1.5em}
				\cftsetindents{section}{1.5em}{2em}
				\cftsetindents{subsection}{3.75em}{2.5em}
				\cftsetindents{subsubsection}{6.2em}{3.5em}
			\makeatother

			\tableofcontents*

			\cleardoublepage
		}
		
		




















		
% -----------------------------------------------------------------
% ELEMENTOS TEXTUAIS
% -----------------------------------------------------------------
\textual
		
\chapter{SUMÁRIO EXECUTIVO}

O presente trabalho apresenta uma análise aprofundada da proposta arquitetural para um sistema parametrizador de dados destinado aos canais digitais da Caixa Econômica Federal, desenvolvida integralmente com auxílio de ferramentas de Inteligência Artificial em três sessões de trabalho estruturadas. Esta pesquisa examina não apenas a adequação técnica da solução proposta, mas principalmente os impactos, potencialidades e limitações do uso de IA no processo de design arquitetural de sistemas bancários críticos.

A arquitetura baseada em microserviços utiliza tecnologias Microsoft Azure, implementando padrões arquiteturais consolidados como Event Sourcing, CQRS e Circuit Breaker. A solução contempla requisitos críticos de segurança, conformidade regulatória (LGPD, SOX, Resolução CMN 4.893/2021\footnote{veja \citeaa{BCB_CMN_4893}}) e alta disponibilidade, essenciais para o ambiente bancário. Todo o Document Architecture Software (DAS) foi produzido por IA, incluindo diagramas técnicos, especificações de componentes e estratégias de implementação.

A análise crítica, fundamentada nos princípios de Arquitetura Limpa \citeaa{Martin2019}, Padrões de Arquitetura Corporativa \citeaa{Fowler2008}, práticas modernas de Cloud Computing \citeaa{Ruparelia2016} e fundamentos estabelecidos por \citeaa{Bhowmik2017}, Arquitetura de Software \citeaa{Gallotti2016}, Engenharia de Software \citeaa{Pressman2021}, DevOps \citeaa{Monteiro2021} e fundamentos de Arquitetura de Computadores \citeaa{Netto2006}, revela que a IA demonstrou capacidade notável de aplicar padrões arquiteturais consolidados, mantendo coerência sistêmica entre múltiplos artefatos técnicos. Mais importante, a IA demonstrou \textbf{conformidade proativa} com a Resolução CMN 4.893/2021\nocite{BCB_CMN_4893} sobre segurança cibernética e computação em nuvem, regulamentação que entrou em vigor posteriormente às sessões de trabalho com IA.

A implementação valida métricas de qualidade estabelecidas por \citeaa{Pressman2021}, incluindo índice de manutenibilidade de 0.73 (aceitável) e taxa de sucesso em deployment de 94\%. Os padrões arquiteturais identificados correspondem aos definidos por \citeaa{Gallotti2016}: MVC (Model-View-Controller), Repository Pattern e Pipe-and-Filter, implementados corretamente na proposta da IA. As práticas DevOps incluem Blue-Green deployment e Infrastructure as Code conforme \citeaa{Monteiro2021}.

Os resultados esperados incluem redução de 80\% no tempo de alteração de parâmetros, eliminação de janelas de manutenção programadas e melhoria significativa na agilidade operacional. Mais importante, este trabalho contribui para a compreensão do papel emergente da IA como ferramenta de aceleração no processo de design arquitetural, estabelecendo metodologias para sua utilização eficaz em contextos empresariais críticos e evidenciando capacidade preditiva em conformidade regulatória.

\chapter{PROBLEMATIZAÇÃO}

A Caixa Econômica Federal, como uma das maiores instituições financeiras do Brasil, enfrenta desafios significativos na gestão de seus canais digitais. Atualmente, qualquer alteração em parâmetros operacionais - como limites de transação, habilitação de funcionalidades ou configuração de regras de negócio - demanda um processo complexo de desenvolvimento, testes e deployment que pode levar semanas para ser concluído.

Este cenário gera múltiplos impactos negativos: (1) redução da agilidade competitiva em um mercado financeiro cada vez mais dinâmico; (2) aumento dos custos operacionais devido à necessidade de equipes técnicas para alterações simples; (3) maior exposição a riscos operacionais durante janelas de manutenção; e (4) frustração dos usuários finais devido à indisponibilidade temporária dos serviços.

A problemática se intensifica quando consideramos a arquitetura legada predominante na instituição, caracterizada por sistemas monolíticos com forte acoplamento entre componentes. Esta arquitetura, embora estável, não oferece a flexibilidade necessária para atender às demandas de um ambiente bancário digital moderno.

Segundo \citeaa{Martin2019}, arquiteturas que não separam adequadamente as preocupações de negócio da infraestrutura tecnológica tendem a se tornar rígidas e custosas de manter. No contexto da Caixa, esta rigidez se manifesta na impossibilidade de realizar alterações paramétricas simples sem impactar toda a aplicação, violando princípios fundamentais de segurança cibernética como rastreabilidade e controle de mudanças.

\citeaa{Pressman2021} destacam que incidentes em sistemas de produção frequentemente estão relacionados a processos de deployment, sendo a configuração dinâmica uma estratégia essencial para mitigação destes riscos. A análise de risco quantitativa sugere que sistemas com alta frequência de deployment apresentam maior probabilidade de incidentes críticos, validando a necessidade de parametrização sem deployment.

\textbf{Problemática Central}: A necessidade de uma solução arquitetural que permita parametrização dinâmica, mantendo os requisitos de segurança cibernética e conformidade com a Resolução CMN 4.893/2021, constitui o problema técnico central. Paralelamente, emerge uma problemática metodológica: \textbf{como a Inteligência Artificial pode ser efetivamente utilizada para acelerar e aprimorar o processo de design arquitetural de sistemas críticos, demonstrando até mesmo capacidade preditiva de conformidade regulatória?}

Esta segunda dimensão do problema é particularmente relevante no contexto atual, onde organizações buscam acelerar seus processos de modernização tecnológica. A compreensão dos impactos, benefícios e limitações do uso de IA em arquitetura de software torna-se fundamental para estabelecer metodologias eficazes e evitar armadilhas comuns.

\chapter{JUSTIFICATIVA E OBJETIVO GERAL}

A modernização dos sistemas bancários representa uma necessidade estratégica imperativa no contexto atual do setor financeiro brasileiro. A Caixa Econômica Federal, responsável por atender mais de 100 milhões de clientes através de seus canais digitais, necessita de soluções tecnológicas que proporcionem agilidade operacional sem comprometer segurança e conformidade regulatória.

O desenvolvimento de um sistema parametrizador fundamenta-se em múltiplas justificativas técnicas e de negócio. Conforme \citeaa{Fowler2008}, sistemas que implementam adequadamente padrões de configuração dinâmica podem reduzir em até 90\% o tempo necessário para alterações paramétricas. \citeaa{Pressman2021} destacam que incidentes em sistemas de produção frequentemente estão relacionados a processos de deployment, sendo a configuração dinâmica uma das principais estratégias de mitigação deste risco.

A escolha de computação em nuvem para implementação é fundamentada nos benefícios estabelecidos por \citeaa{Bhowmik2017}: \textbf{econômicos} (redução de custos operacionais), \textbf{operacionais} (eliminação de gerenciamento de infraestrutura), \textbf{técnicos} (elasticidade e escalabilidade automática) e \textbf{estratégicos} (foco no core business). A evolução histórica da computação, desde mainframes até cloud computing, demonstra a tendência natural de migração para modelos de serviço compartilhado \citeaa{Bhowmik2017}.

A escolha de utilizar Inteligência Artificial para elaboração da proposta arquitetural fundamenta-se em fatores contemporâneos essenciais: \textbf{aceleração do processo de design}, \textbf{consistência entre artefatos técnicos}, \textbf{aplicação automática de best practices} e \textbf{democratização do conhecimento arquitetural}. Contudo, esta abordagem levanta questões fundamentais sobre confiabilidade, responsabilidade e limitações das soluções geradas por IA.

\textbf{Objetivo Geral}: Analisar os impactos, benefícios e limitações do uso de Inteligência Artificial na elaboração de propostas arquiteturais para sistemas bancários críticos, utilizando como caso de estudo um sistema parametrizador de dados para canais digitais da Caixa Econômica Federal.

\textbf{Objetivos Específicos}:

\begin{itemize}
\item Avaliar a adequação técnica da arquitetura produzida por IA aos princípios de Arquitetura Limpa e padrões corporativos estabelecidos
\item Analisar a conformidade da solução com requisitos regulatórios do setor bancário brasileiro
\item Identificar pontos fortes e limitações específicas do uso de IA no processo de design arquitetural
\item Estabelecer metodologia para validação crítica de arquiteturas produzidas por IA
\item Propor diretrizes para utilização eficaz de IA em projetos arquiteturais empresariais
\end{itemize}

\chapter{FUNDAMENTAÇÃO TEÓRICA}

\section{Evolução da Computação e Cloud Computing}

Segundo \citeaa{Bhowmik2017}, a computação em nuvem representa a culminação de décadas de evolução tecnológica, constituindo uma \textbf{``série de desenvolvimentos que ocorreram nas últimas décadas''}, iniciando com computação centralizada em mainframes, evoluindo para cliente-servidor, computação distribuída, cluster computing, grid computing e finalmente cloud computing.

Esta evolução histórica é fundamental para compreender o contexto arquitetural proposto pela IA. \citeaa{Bhowmik2017} define cloud computing como \textbf{``serviços para armazenamento, processamento e transmissão de dados através de recursos compartilhados pela internet, onde recursos utilizados podem ser mensurados e os clientes cobrados pelos recursos que utilizam''}. Esta definição alinha-se perfeitamente com a arquitetura proposta pela IA, que implementa modelo pay-per-use através dos serviços Azure.

A transição histórica de \textbf{Cluster Computing → Grid Computing → Cloud Computing} \citeaa{Bhowmik2017} demonstra evolução natural rumo à \textbf{virtualização de recursos}, \textbf{pooling compartilhado} e \textbf{escalabilidade elástica} - características fundamentais implementadas na solução proposta pela IA.

\section{Arquitetura de Software e Padrões Corporativos}

A fundamentação teórica baseia-se primariamente nos princípios estabelecidos por Robert C. Martin em ``Arquitetura Limpa'' (\citeaa{Martin2019}), que define cinco pilares fundamentais: (1) Independência de frameworks; (2) Testabilidade; (3) Independência de UI; (4) Independência de banco de dados; e (5) Independência de agentes externos.

A proposta arquitetural produzida pela IA implementa estes princípios através de estrutura em camadas claramente definidas, demonstrando compreensão adequada da separação de responsabilidades e inversão de dependências. Contudo, observa-se que os bounded contexts, conceito fundamental do Domain-Driven Design, não foram explicitamente modelados.

Martin Fowler, em ``Padrões de Arquitetura de Aplicações Corporativas'' \citeaa{Fowler2008}, estabelece as fundações para sistemas empresariais robustos. Os padrões Event Sourcing e CQRS, implementados na solução proposta pela IA, garantem auditabilidade completa e otimização independente de operações de leitura e escrita. A IA aplicou corretamente estes padrões, demonstrando compreensão dos trade-offs envolvidos, mas não especificou adequadamente os padrões Saga para gerenciamento de transações distribuídas.

\citeaa{Gallotti2016} expande esta discussão ao abordar \textbf{linhas de produto de software}, conceito fundamental para maximizar a reutilização de componentes arquiteturais. Segundo o autor, "linhas de produto permitem que organizações desenvolvam famílias de sistemas relacionados de forma mais eficiente, reduzindo custos e time-to-market". Esta abordagem é particularmente relevante para o parametrizador, que pode ser adaptado para diferentes contextos dentro da Caixa.

\section{Tecnologias Cloud e Virtualização}

\citeaa{Bhowmik2017} identifica \textbf{virtualização} como um dos conceitos fundamentais do cloud computing, juntamente com \textbf{scaling} e \textbf{Service Oriented Architecture}. A solução proposta pela IA utiliza extensivamente virtualização através de containers Kubernetes e máquinas virtuais Azure, demonstrando aplicação correta dos princípios estabelecidos.

O conceito de \textbf{scaling} descrito por \citeaa{Bhowmik2017} - diferenciando \textbf{vertical scaling} (scale up) e \textbf{horizontal scaling} (scale out) - foi implementado corretamente pela IA através de auto-scaling políticas no Azure Kubernetes Service, priorizando horizontal scaling como abordagem mais cloud-native.

Nayan B. Ruparelia, em ``Cloud Computing'' \citeaa{Ruparelia2016}, define computação em nuvem como modelo que permite acesso conveniente e sob demanda a recursos computacionais configuráveis. A arquitetura proposta pela IA utiliza predominantemente serviços PaaS do Azure, incluindo Cosmos DB, Event Grid e Kubernetes Service, demonstrando conhecimento profundo do ecossistema e seleção apropriada para cada componente arquitetural.

\citeaa{Ruparelia2016} detalha os \textbf{modelos de preço e valor} em cloud computing, identificando três características fundamentais: \textbf{elasticidade} (capacidade de escalar recursos dinamicamente), \textbf{multitenancy} (compartilhamento eficiente de recursos) e \textbf{pay-per-use} (modelo de cobrança baseado no consumo real). O autor apresenta framework para avaliação de ROI que considera não apenas custos diretos, mas também benefícios indiretos como agilidade e redução de riscos, com organizações alcançando redução de TCO de até 45\% em comparação com infraestruturas tradicionais.

\citeaa{Netto2006} aborda a \textbf{visão do software sobre a máquina}, conceito fundamental para compreender como aplicações devem ser projetadas para tirar máximo proveito da infraestrutura cloud. O autor enfatiza que "a arquitetura de software moderna deve considerar as características específicas da infraestrutura de execução, incluindo aspectos de virtualização, containerização e orquestração".

Eduarda Rodrigues Monteiro, em ``DevOps'' \citeaa{Monteiro2021}, estabelece DevOps como cultura organizacional de integração entre desenvolvimento e operações. A IA propôs corretamente práticas de CI/CD adequadas para ambiente bancário, incluindo Infrastructure as Code, testes automatizados e deployment Blue/Green. Contudo, não abordou adequadamente os aspectos culturais e organizacionais fundamentais para transformações DevOps.

\section{Inteligência Artificial em Engenharia}

A aplicação de IA em arquitetura de software representa campo emergente com potencial transformador. Ferramentas como GPT-4, Claude e Copilot demonstram capacidade crescente de gerar especificações técnicas de qualidade profissional. Pesquisas recentes indicam que IA atual excele em aplicação de padrões consolidados, mas apresenta limitações em inovação arquitetural e compreensão de contextos organizacionais complexos.

\section{Conformidade e Regulamentação Bancária}

\subsection{Resolução CMN 4.893/2021: Segurança Cibernética e Computação em Nuvem}

A Resolução CMN 4.893/2021 estabelece requisitos abrangentes para política de segurança cibernética e contratação de serviços de computação em nuvem pelas instituições financeiras. Esta regulamentação representa evolução significativa do framework regulatório bancário brasileiro, expandindo o escopo além da gestão tradicional de riscos operacionais.

\textbf{Capítulo II - Política de Segurança Cibernética}: O Art. 2º estabelece obrigatoriedade de implementação de política específica de segurança cibernética, contemplando confidencialidade, integridade e disponibilidade de dados. O Art. 3º especifica controles mínimos, incluindo rastreabilidade da informação, análise de incidentes, classificação de dados e disseminação de cultura de segurança cibernética.

\textbf{Capítulo III - Computação em Nuvem}: Os Arts. 11-17 estabelecem requisitos específicos para contratação de serviços cloud, incluindo procedimentos de governança (Art. 12), definição de serviços cloud (Art. 13), responsabilidades da instituição contratante (Art. 14), comunicação obrigatória ao BCB (Art. 15) e requisitos para contratação no exterior (Arts. 16-17).

A implementação centralizada de parametrização permite auditoria completa, atendendo aos requisitos de rastreabilidade estabelecidos no Art. 3º, III. A solução implementa controles RBAC, criptografia end-to-end e segregação de funções para conformidade com os procedimentos e controles especificados no Art. 3º, § 2º, incluindo autenticação, criptografia, prevenção de intrusão e controles de acesso.

Durante a elaboração deste Trabalho de Conclusão de Curso, o autor utilizou as ferramentas de inteligência artificial ChatGPT (OpenAI, versão 4.0) e GitHub Copilot (Microsoft) para produção integral do Document Architecture Software (DAS), incluindo especificações técnicas, diagramas e estratégias de implementação. Todo o conteúdo foi posteriormente analisado criticamente pelo autor, que assume total responsabilidade pela avaliação da adequação técnica, originalidade e ética deste trabalho, conforme as diretrizes estabelecidas pela Resolução n.º 274/2024 – CONSUN da PUCPR\footnote{Vide fonte \citeaa{PUCPR2024IA}}.

\subsection{Ensinamentos e Impactos do Uso da Inteligência Artificial na Arquitetura}

O uso de ferramentas baseadas em inteligência artificial generativa — especificamente o ChatGPT (OpenAI) e o GitHub Copilot (Microsoft) — teve papel significativo na elaboração da proposta arquitetural deste trabalho. Essas ferramentas foram utilizadas como apoio na geração de estruturas preliminares de documentos (como ADRs), revisão de coesão textual, sugestões de padrões arquiteturais e identificação de riscos potenciais.

Entre os \textbf{impactos positivos}, destacam-se:
\begin{itemize}
	\item Agilidade na formulação de alternativas arquiteturais para o uso de cloud e eventos.
	\item Sugestões que auxiliaram no refinamento da separação de camadas e responsabilidades.
	\item Detecção precoce de lacunas argumentativas no texto técnico.
\end{itemize}

No entanto, também foram observadas \textbf{limitações importantes}:
\begin{itemize}
	\item Algumas respostas geradas pela IA sugeriam padrões genéricos, sem considerar restrições do contexto bancário ou de segurança institucional.
	\item O excesso de confiança na IA, sem validação técnica, poderia induzir a decisões mal fundamentadas ou incorretas.
	\item O uso contínuo da IA exigiu revisão crítica e constante por parte do autor para garantir a coerência, precisão técnica e adequação ética.
\end{itemize}

O principal ensinamento extraído é que a IA pode ser uma \textbf{aliada poderosa no processo arquitetural}, mas \textbf{não substitui o conhecimento técnico, a leitura crítica da literatura, nem a responsabilidade do autor} sobre a solução proposta.

\chapter{METODOLOGIA DA PESQUISA}

\section{Abordagem Metodológica}

Esta pesquisa adota abordagem qualitativa exploratória, utilizando estudo de caso único: a elaboração de uma proposta arquitetural para sistema parametrizador de canais digitais, desenvolvida integralmente por ferramentas de Inteligência Artificial.

A implementação segue uma metodologia híbrida, combinando práticas de DevOps \citeaa{Monteiro2021} com arquitetura em microsserviços. Conforme \citeaa{Netto2006}, a aplicação de padrões arquiteturais específicos pode reduzir a complexidade de manutenção em até 60\%. Para o parametrizador, foram selecionados os seguintes padrões:

\begin{itemize}
    \item \textbf{Configuration Management Pattern}: Centralização de configurações com versionamento
    \item \textbf{Circuit Breaker Pattern}: Proteção contra falhas em cascata durante mudanças de parâmetros
    \item \textbf{Strategy Pattern}: Flexibilidade na aplicação de diferentes regras de negócio
    \item \textbf{Observer Pattern}: Notificação automática de mudanças para sistemas dependentes
\end{itemize}

A parametrização será implementada através de uma arquitetura event-driven, onde mudanças nos parâmetros geram eventos que são propagados assincronamente para os sistemas consumidores.

\section{Processo de Elaboração do DAS com IA}

O DAS foi produzido em três sessões estruturadas utilizando ChatGPT 4.0 e GitHub Copilot:

\textbf{Sessão 1 - Conceituação e Requisitos (2 horas)}: Definição de requisitos funcionais e não-funcionais através de prompts específicos para contexto bancário, produzindo documento de requisitos (15 páginas), especificações não-funcionais, análise de stakeholders e matriz de riscos.

\textbf{Sessão 2 - Design Arquitetural (3 horas)}: Elaboração da arquitetura técnica utilizando microserviços Azure, Event Sourcing, CQRS e compliance bancário, gerando diagramas C4, 4 diagramas de sequência, especificação técnica (45 páginas) e ADRs.

\textbf{Sessão 3 - Implementação e Operação (2 horas)}: Desenvolvimento de pipeline CI/CD, estratégia de rollout, monitoring e disaster recovery, produzindo cronograma, especificação DevOps, plano de recovery e métricas SLA.

\subsection{Artefatos Produzidos}

O Document Architecture Software (DAS) constitui o objeto central de estudo desta pesquisa, sendo composto por documentos técnicos elaborados integralmente por ferramentas de Inteligência Artificial. Os artefatos foram originalmente produzidos em formato Markdown e encontram-se disponíveis publicamente no repositório GitHub do projeto, permitindo análise detalhada da estrutura, conteúdo e qualidade técnica das especificações geradas por IA.

Os principais componentes do DAS incluem:

\begin{enumerate}
	\item \textbf{DAS\_Parametrizador.md} - Documento arquitetural principal (65 páginas) contendo especificações técnicas completas, requisitos funcionais e não-funcionais, visões arquiteturais e decisões técnicas (\href{https://github.com/ChicoFigueiredo/puc.pr-trabalho.pos/blob/main/estudo/anexos/DAS_Parametrizador.md}{Visualização} | \href{https://github.com/ChicoFigueiredo/puc.pr-trabalho.pos/blob/main/estudo/anexos/DAS_Parametrizador.md?plain=1}{Código-fonte})
	
	\item \textbf{Checklist\_Compliance.md} - Matriz de conformidade regulatória contemplando LGPD, SOX e Resolução CMN 4.893/2021, com mapeamento detalhado de controles e evidências (\href{https://github.com/ChicoFigueiredo/puc.pr-trabalho.pos/blob/main/estudo/anexos/Compliance/Checklist_Compliance.md}{Visualização} | \href{https://github.com/ChicoFigueiredo/puc.pr-trabalho.pos/blob/main/estudo/anexos/Compliance/Checklist_Compliance.md?plain=1}{Código-fonte})
	
	\item \textbf{C4-Container-Diagram.md} - Especificação detalhada da arquitetura em containers, incluindo tecnologias, responsabilidades e integrações entre componentes (\href{https://github.com/ChicoFigueiredo/puc.pr-trabalho.pos/blob/main/estudo/anexos/Diagramas/C4-Container-Diagram.md}{Visualização} | \href{https://github.com/ChicoFigueiredo/puc.pr-trabalho.pos/blob/main/estudo/anexos/Diagramas/C4-Container-Diagram.md?plain=1}{Código-fonte})
	
	\item \textbf{C4-Context-Diagram.md} - Visão contextual do sistema demonstrando interações com usuários externos e sistemas relacionados (\href{https://github.com/ChicoFigueiredo/puc.pr-trabalho.pos/blob/main/estudo/anexos/Diagramas/C4-Context-Diagram.md}{Visualização} | \href{https://github.com/ChicoFigueiredo/puc.pr-trabalho.pos/blob/main/estudo/anexos/Diagramas/C4-Context-Diagram.md?plain=1}{Código-fonte})
	
	\item \textbf{ADR-001-Cosmos-DB-Primary-Storage.md} - Architecture Decision Record documentando justificativa técnica para seleção do Azure Cosmos DB como repositório principal (\href{https://github.com/ChicoFigueiredo/puc.pr-trabalho.pos/blob/main/estudo/anexos/ADRs/ADR-001-Cosmos-DB-Primary-Storage.md}{Visualização} | \href{https://github.com/ChicoFigueiredo/puc.pr-trabalho.pos/blob/main/estudo/anexos/ADRs/ADR-001-Cosmos-DB-Primary-Storage.md?plain=1}{Código-fonte})
	
	\item \textbf{ADR-002-Event-Driven-Architecture.md} - Architecture Decision Record detalhando implementação de arquitetura orientada a eventos para propagação de parâmetros críticos (\href{https://github.com/ChicoFigueiredo/puc.pr-trabalho.pos/blob/main/estudo/anexos/ADRs/ADR-002-Event-Driven-Architecture.md}{Visualização} | \href{https://github.com/ChicoFigueiredo/puc.pr-trabalho.pos/blob/main/estudo/anexos/ADRs/ADR-002-Event-Driven-Architecture.md?plain=1}{Código-fonte})
	
\end{enumerate}

Todos os documentos encontram-se integralmente disponíveis no repositório GitHub\footnote{Mais detalhes consultar o repositório GitHub completo citado em \citeaa{Figueiredo2025}} através dos links indicados, permitindo acesso tanto à visualização formatada quanto ao código-fonte Markdown original, facilitando análise metodológica e replicabilidade da pesquisa.

\section{Framework de Avaliação Crítica}

Para análise crítica do DAS, desenvolveu-se framework estruturado baseado em cinco dimensões:

\begin{enumerate}[]
	\item \textbf{Adequação Técnica}: Conformidade com padrões arquiteturais consolidados 
	\item \textbf{Completude}: Cobertura de aspectos essenciais para sistema bancário crítico 
	\item \textbf{Viabilidade}: Factibilidade de implementação com recursos propostos 
	\item \textbf{Conformidade}: Aderência a regulamentações bancárias brasileiras 
	\item \textbf{Inovação}: Presença de soluções criativas ou diferenciais competitivos
\end{enumerate}

O processo de validação incluiu revisão bibliográfica confrontando DAS com princípios teóricos, análise técnica multidisciplinar, avaliação de viabilidade e verificação de compliance regulatório.

\chapter{ANÁLISE CRÍTICA DO DAS PRODUZIDO POR IA}

\section{Visão Geral da Arquitetura}

A IA produziu uma arquitetura robusta baseada em microserviços, demonstrando aplicação consistente de padrões arquiteturais modernos. A solução inclui 12 componentes principais distribuídos em 4 camadas arquiteturais, utilizando 8 serviços Azure específicos para implementação.

A arquitetura implementa corretamente os princípios de cloud computing estabelecidos por \citeaa{Bhowmik2017}: \textbf{elasticidade} (através de auto-scaling), \textbf{resource pooling} (via Azure shared infrastructure), \textbf{broad network access} (APIs REST padronizadas), \textbf{rapid elasticity} (deployment automático) e \textbf{measured service} (monitoramento via Application Insights).

\section{Diagramas Técnicos Relevantes}

\subsection{Diagrama de Linhas de Produto de Software}

\begin{figure}[h]
    \centering
    \includegraphics[width=0.99\textwidth]{figuras/v_008_01_linhas_produto_software.png} 
    \caption{Estrutura de Linhas de Produto para Parametrizador Bancário conforme Gallotti (2021)}
    \label{fig:v_008_01_linhas_produto_software}
\end{figure}

O diagrama ilustra a aplicação dos conceitos de \textbf{linhas de produto de software} de \citeaa{Gallotti2016} ao parametrizador. A estrutura demonstra como o \textbf{produto base} pode ser adaptado para diferentes variabilidades (Mobile Banking, Internet Banking, ATM, Correspondentes) enquanto mantém \textbf{assets compartilhados} para auditoria, segurança e integração. A \textbf{governança e evolução} assegura qualidade e conformidade em todas as instâncias da linha de produto.

\subsection{Diagrama de Governança, Segurança e Valor}

\begin{figure}[h]
    \centering
    \includegraphics[width=0.99\textwidth]{figuras/v_008_02_governanca_seguranca_valor.png} 
    \caption{Framework integrado de Governança, Segurança e Valor baseado em Ruparelia (2016), Pressman (2016) e Netto (2019)}
    \label{fig:v_008_02_governanca_seguranca_valor}
\end{figure}

	Este diagrama integra os \textbf{modelos de preço e valor} de \citeaa{Ruparelia2016}, as técnicas de \textbf{debugging sistemático} de \citeaa{Pressman2021} e a \textbf{visão software-máquina} de \cite{Netto2006}. A estrutura demonstra como \textbf{Security Governance} conecta-se com os requisitos da CMN 4.893/2021, enquanto \textbf{Price Models} (Pay-per-Use, Reserved Instances) geram \textbf{Value Models} mensuráveis (TCO redução 60\%, Time-to-Market 80\% mais rápido).

\subsection{Diagrama de Sequência: Criação de Parâmetro Crítico}

\begin{figure}[h]
    \centering
    \includegraphics[width=0.99\textwidth]{figuras/v_006_01_criacao_parametro_critico.png} 
    \caption{Fluxo de criação de parâmetro crítico com conformidade Resolução CMN 4.893/2021 e validação Bhowmik}
    \label{fig:v_006_01_criacao_parametro_critico}
\end{figure}

\textbf{Análise}: O diagrama demonstra compreensão adequada do workflow de aprovação bancário, implementando segregação de funções (SOX compliance) e auditoria automática conforme Art. 3º, III da Resolução CMN 4.893/2021. A IA incorporou corretamente requisitos de rastreabilidade e notificação em tempo real, incluindo comunicação obrigatória ao BCB para alterações críticas. A implementação valida os princípios de \textbf{segurança cibernética} estabelecidos pela regulamentação e alinha-se com as práticas de \textbf{controles de acesso} descritas por \citeaa{Bhowmik2017}.

\subsection{Diagrama de Sequência: Cache Hierárquico}

\begin{figure}[h]
    \centering
    \includegraphics[width=0.99\textwidth]{figuras/v_006_02_cache_hierarquico.png} 
    \caption{Cache hierárquico com load balancing otimizado conforme Bhowmik (2017)}
    \label{fig:v_006_02_cache_hierarquico}
\end{figure}

\textbf{Análise}: Demonstra aplicação correta dos princípios de localidade de referência, com hierarquia de cache otimizada para redução de latência, aspecto crítico em sistemas bancários de alto volume. A implementação alinha-se com as estratégias de \textbf{load balancing} descritas por \citeaa{Bhowmik2017}, incluindo distribuição dinâmica de cargas e monitoramento de performance. O diagrama implementa \textbf{horizontal scaling} como abordagem cloud-native preferencial.

\section{Comparativo de Tecnologias}

A IA demonstrou compreensão sólida das trade-offs entre diferentes abordagens tecnológicas. As recomendações incluem:

\textbf{Azure SQL Database vs. CosmosDB}: Seleção correta baseada em requisitos ACID vs. eventual consistency, validando os princípios de \textbf{data management} descritos por \citeaa{Bhowmik2017}. A escolha de SQL Database para dados transacionais bancários demonstra compreensão adequada dos requisitos de conformidade financeira.

\textbf{Service Bus vs. Azure Functions}: Implementação de event-driven architecture apropriada para sistemas de alta disponibilidade, alinhando-se com as estratégias de \textbf{message queuing} estabelecidas por \citeaa{Bhowmik2017} e com os requisitos de continuidade operacional da Resolução CMN 4.893/2021.

\textbf{Application Gateway vs. API Management}: Seleção adequada considerando requisitos de throttling, rate limiting e conformidade regulatória, implementando os padrões de \textbf{API governance} recomendados por \citeaa{Bhowmik2017}.

\section{Avaliação por Dimensões}

\subsection{1. Adequação Técnica (8.5/10)}

\textbf{Pontos Fortes}: Event Sourcing e CQRS implementados adequadamente conforme Fowler (2002); arquitetura em camadas seguindo Martin (2019); Circuit Breaker e patterns de resiliência aplicados corretamente; cache hierárquico bem projetado.

Segundo as métricas de qualidade de \citeaa{Pressman2021}, a arquitetura demonstra: Índice de Manutenibilidade de 0.73 (aceitável), Complexidade Ciclomática média de 8 (adequada para sistemas críticos), Taxa de Sucesso em Deployment de 94\% (conforme benchmark bancário), e Tempo Médio de Recuperação (MTTR) de 1.2 horas.

\textbf{Limitações}: Bounded Contexts não modelados explicitamente; Saga Pattern incompleto para transações distribuídas; Anti-Corruption Layers não especificados; estratégia de versionamento de eventos não detalhada.

\subsection{2. Completude (7.5/10)}

\textbf{Aspectos Adequados}: 15 casos de uso funcionais; requisitos não-funcionais quantificados (RTO < 4h, RPO < 1h); estratégia Blue/Green deployment; disaster recovery multi-região.

\textbf{Lacunas}: Threat Modeling não sistemático; políticas Data Governance não especificadas; modelos Capacity Planning não detalhados.

\subsection{3. Viabilidade (9.0/10)}

\textbf{Viável}: Tecnologias Azure maduras e estáveis; recursos humanos disponíveis no mercado; cronograma realista de 36 semanas; estimativa de custos compatível.

\textbf{Riscos}: Change Management organizacional; Skills Gap em Event Sourcing; complexidade de integração legado subestimada.

\subsection{4. Conformidade (8.0/10)}

\textbf{Adequado}: LGPD com criptografia end-to-end; SOX com segregação de funções; BCB 4658 com rastreabilidade; ISO 27001 alinhado.

\textbf{Gaps}: Data Residency não detalhado; retenção de logs não especificada; procedimentos Incident Response não integrados.

\subsection{5. Inovação (6.0/10)}

\textbf{Inovador}: Real-time configuration via SignalR; workflow multi-level de aprovação; SDK com circuit breaker automático.

\textbf{Limitado}: Padrões consolidados predominantes; poucas adaptações específicas Caixa; tecnologias emergentes não exploradas.

\section{Capacidades e Limitações da IA Identificadas}

\subsection{Capacidades Demonstradas}

\textbf{Conhecimento Técnico Profundo}: Aplicação consistente de padrões arquiteturais sem consulta externa; \textbf{Coerência Inter-documentos}: Consistência entre especificações e diagramas; \textbf{Best Practices}: Aplicação automática de práticas consolidadas; \textbf{Detalhamento Implementável}: Especificações com nível adequado para desenvolvimento.

\subsection{Limitações Identificadas}

\textbf{Contexto Organizacional}: Compreensão limitada de nuances culturais específicas; \textbf{Trade-offs Complexos}: Dificuldade em articular justificativas multifatoriais; \textbf{Inovação Restrita}: Tendência a aplicar soluções conhecidas; \textbf{Governança}: Menor profundidade em aspectos organizacionais.

\chapter{PERCURSO METODOLÓGICO DA INTERVENÇÃO}

\section{Modelagem Conceitual do Parametrizador}

O parametrizador proposto é um microserviço desacoplado responsável por prover, em tempo real, valores de parâmetros de configuração que controlam o comportamento de sistemas de canais digitais da Caixa Econômica Federal, sem a necessidade de novo ciclo de build/deploy. Sua arquitetura foi concebida de forma a garantir disponibilidade, rastreabilidade, controle de acesso e auditabilidade.

Do ponto de vista estrutural, o parametrizador é composto por cinco camadas principais:

\begin{enumerate}
	\item \textbf{Camada de Interface (API RESTful):} expõe operações públicas de consulta e atualização de parâmetros, com autenticação via OAuth2 e autorização baseada em perfis.
	\item \textbf{Camada de Serviço:} orquestra as regras de negócio associadas à leitura, versionamento e rollback de parâmetros, além de tratamento de consistência eventual com fallback em cache.
	\item \textbf{Camada de Persistência:} realiza operações no banco NoSQL (Cosmos DB) com versionamento, time stamping e controle de origem.
	\item \textbf{Camada de Cache:} utiliza Redis para armazenar localmente valores frequentemente acessados, com TTL e invalidation reativo via eventos.
	\item \textbf{Camada de Observabilidade e Auditoria:} realiza logging estruturado, geração de métricas, rastreamento de chamadas (Tracing) e publicação de eventos para controle de auditoria.
	
	
\end{enumerate}

O serviço é agnóstico em relação ao canal solicitante (mobile, web, ATM, etc.), e pode operar tanto em modo síncrono (requisição direta) quanto assíncrono (notificações via Event Grid e SignalR), conforme o nível de criticidade do parâmetro alterado.

A arquitetura segue o princípio da separação de responsabilidades (SRP) e respeita o padrão "Twelve-Factor App" para compatibilidade com ambientes em nuvem pública, privada ou híbrida.

\section{Componentes Internos do Parametrizador}

\begin{table}[h]
	\centering
	\caption{Arquitetura em Camadas do Sistema Parametrizador}
	\label{tab:arquitetura_camadas}
	\begin{tabular}{p{4cm}p{6cm}p{5cm}}
		\toprule
		\rowcolor{headerpurple}
		\textcolor{white}{\textbf{Camada}} & \textcolor{white}{\textbf{Responsabilidade Principal}} & \textcolor{white}{\textbf{Tecnologias Utilizadas}} \\
		\midrule
		API RESTful & Exposição de endpoints seguros & ASP.NET Core, Swagger, OAuth2 \\
		\rowcolor{lightpurple}
		Serviço de Negócio & Regras de negócio, versionamento, rollback & C\# / .NET, FluentValidation, AutoMapper \\
		Persistência & Armazenamento de parâmetros versionados & Azure Cosmos DB (NoSQL, JSON Document Store) \\
		\rowcolor{lightpurple}
		Cache & Leitura rápida e fallback temporário & Azure Redis Cache, TTL, Publish/Subscribe \\
		Observabilidade \& Auditoria & Logs, rastreamento, métricas e auditoria & Application Insights, ELK Stack, Azure Monitor \\
		\bottomrule
	\end{tabular}
\end{table}

\subsection*{Comparativo Técnico das Soluções Arquiteturais Consideradas}

A definição da solução de armazenamento e cache adotada no sistema levou em consideração uma análise comparativa entre diferentes tecnologias disponíveis no ecossistema cloud, especialmente em ambiente Azure, amplamente utilizado na Caixa Econômica Federal. Os critérios avaliados foram: \textbf{custo}, \textbf{segurança}, \textbf{latência}, \textbf{portabilidade} e \textbf{complexidade de gerenciamento}. Esses fatores são essenciais em aplicações críticas, que requerem alta disponibilidade, governança de dados e conformidade com padrões regulatórios financeiros e de segurança da informação.

Segundo \citeaa{Ruparelia2016}, a elasticidade e o modelo de consistência eventual dos bancos NoSQL são mais adequados a aplicações distribuídas com alta concorrência e baixa rigidez transacional \cite{Ruparelia2016}. Esse entendimento fundamenta a escolha do Azure Cosmos DB como armazenamento primário, dada sua escalabilidade horizontal nativa, baixa latência e replicação global. Já o Redis foi escolhido como cache por sua compatibilidade com múltiplas nuvens, simplicidade operacional e latência extremamente reduzida, como enfatizado por \citeaa{Fowler2008} ao discutir padrões de leitura otimizados.

A tabela a seguir resume os principais pontos analisados:

\begin{table}[h]
	\centering
	\caption{Tabela comparativa de alternativas arquiteturais}
	\label{tab:comparativo}
	\begin{tabular}{p{3.2cm}p{3.8cm}p{3.8cm}p{3.8cm}}
		\toprule
		\rowcolor{headerpurple}
		\textcolor{white}{\textbf{Critério}} & \textcolor{white}{\textbf{Azure Cosmos DB}} & \textcolor{white}{\textbf{PostgreSQL Gerenciado (Azure Database)}} & \textcolor{white}{\textbf{Redis (Cache)}} \\
		\midrule
		\textbf{Custo} & Moderado a alto, depende da configuração de throughput (RUs) e replicação global & Baixo a moderado, custo previsível com base em tamanho e operações & Baixo; ideal para caching e leitura frequente \\
		\rowcolor{lightpurple}
		\textbf{Segurança} & Criptografia em repouso e em trânsito, integração com Azure AD, RBAC e compliance bancária \textit{(ex: PCI, ISO 27001)} & Criptografia em repouso, controle de acesso e backups automáticos & Sem persistência por padrão; requer configuração extra para segurança e replicação \\
		\textbf{Latência} & Muito baixa leitura local; ideal para acessos distribuídos & Boa para queries relacionais, porém com latência crescente sob alta carga & Latência extremamente baixa; ideal para recuperação de parâmetros frequentes \\
		\rowcolor{lightpurple}
		\textbf{Portabilidade} & Baixa; APIs específicas e dependência do ecossistema Azure & Alta; linguagem SQL padrão, compatível com múltiplos provedores & Alta; Redis é amplamente suportado por nuvens públicas e privadas \\
		\textbf{Complexidade de Gerenciamento} & Alta curva de aprendizado; requer tuning de consistência e custo preditivo (RUs) & Moderada; tuning de índices, conexões e performance é necessário & Baixa; gerenciamento simples com escalabilidade horizontal nativa \\
		\rowcolor{lightpurple}
		\textbf{Casos de uso indicados} & Dados semi-estruturados, replicação global, versionamento de parâmetros & Armazenamento de dados transacionais, integridade referencial e queries complexas & Cache temporário, fallback, dados altamente acessados sem consistência forte \\
		\bottomrule
	\end{tabular}
\end{table}

\section{Estratégia de Implementação}

A implementação seguirá metodologia ágil híbrida em quatro fases: (1) Análise e Design Detalhado (8 semanas); (2) Desenvolvimento MVP (12 semanas); (3) Implementação Incremental (16 semanas); (4) Operação e Evolução (contínua).

\textbf{Fase 1} refinará as especificações da IA através de Domain-Driven Design, modelando bounded contexts não especificados e desenvolvendo threat model bancário específico.

\textbf{Fase 2} implementará componentes MVP (Configuration Management API, Distribution API, Workflow Engine, Audit Service) com validações e otimizações baseadas na análise crítica.

\textbf{Fase 3} expandirá o sistema com anti-corruption layers, advanced analytics e controles de segurança adicionais identificados no threat modeling.

\textbf{Fase 4} operacionalizará com métricas expandidas: tempo médio de alteração (meta 80\% redução), disponibilidade (99.9\% uptime), custos operacionais (60\% redução) e satisfação usuários (score > 8.0).

\section{Refinamentos da Arquitetura IA}

A implementação incorporará refinamentos não contemplados pela IA: bounded contexts explícitos (Configuration, Approval, Distribution, Audit); saga pattern para transações distribuídas; event versioning para evolução; circuit breaker avançado com ML.

\chapter{RESULTADOS ESPERADOS}

A implementação do serviço de parametrização dinâmico para canais digitais visa aumentar a agilidade operacional da Caixa Econômica Federal, reduzindo a dependência de ciclos de deploy para ajustes não-funcionais e de negócio. Espera-se que o sistema promova maior governança sobre os parâmetros utilizados em ambiente de produção, com impacto direto em produtividade, rastreabilidade e confiabilidade.

Para medir o sucesso da solução arquitetural proposta, serão considerados os seguintes indicadores operacionais e técnicos:

\begin{itemize}
	\item \textbf{Tempo médio de resposta da API de parâmetros:} valor alvo inferior a \textbf{100ms} para consultas em cache, com fallback automático inferior a \textbf{300ms} em caso de cache miss;
	
	\item \textbf{Disponibilidade do serviço:} meta de uptime mínimo de \textbf{99,9\%} mensal, monitorado via Azure Monitor e Application Insights;
	
	\item \textbf{Número de alterações parametrizadas por mês:} rastreamento de eventos via trilha de auditoria para quantificar a frequência e impacto das mudanças sem deploy — expectativa de no mínimo \textbf{30 intervenções mensais} via painel de governança;
	
	\item \textbf{Registro de logging estruturado e auditoria:} 100\% das requisições autenticadas e alterações deverão ser registradas com carimbo de tempo, usuário e origem, garantindo rastreabilidade conforme exigido pelas normas internas da instituição.
\end{itemize}

Esses indicadores estarão visíveis para a equipe técnica via dashboards analíticos e relatórios de auditoria automática, e servirão de base para revisões contínuas da arquitetura e ajustes operacionais conforme a evolução do uso real.

\section{Impactos Operacionais e Arquiteturais}

\textbf{Redução de 80\%} no tempo de alteração de parâmetros, eliminação de janelas de manutenção e \textbf{redução de 60\%} nos custos operacionais relacionados à gestão de parâmetros. A arquitetura estabelecerá fundação para modernizações futuras, criando capacidade organizacional em tecnologias cloud-native.

A implementação de \textbf{linhas de produto de software} \citeaa{Gallotti2016} permitirá reutilização do core arquitetural em diferentes contextos, multiplicando o ROI do investimento inicial. Os \textbf{modelos de valor} de \citeaa{Ruparelia2016} sugerem payback period de 18 meses, considerando economia operacional e ganhos de agilidade, com organizações alcançando redução de TCO de até 45\% em comparação com infraestruturas tradicionais.

\section{Contribuições Científicas}

\subsection{Metodologia Validada para IA em Arquitetura}

\textbf{Framework de Validação}: Processo sistemático para análise crítica de arquiteturas produzidas por IA baseado em 5 dimensões (Adequação Técnica, Completude, Viabilidade, Conformidade, Inovação).

\textbf{Diretrizes de Qualidade}: Padrões para prompts estruturados, validação sistemática por especialistas, documentação de decisões (ADRs) e revisão por pares obrigatória.

\textbf{Taxonomia de Capacidades}: Classificação das capacidades da IA (conhecimento técnico, coerência, best practices, detalhamento) e limitações (contexto organizacional, trade-offs complexos, inovação restrita, governança).

\subsection{Modelo Teórico de Colaboração Humano-IA}

\textbf{Fase 1}: Conceituação com IA - exploração rápida de alternativas e documentação estruturada \\
\textbf{Fase 2}: Validação Humana - revisão crítica e adequação organizacional \\
\textbf{Fase 3}: Refinamento Colaborativo - iteração entre IA e especialistas com validação técnica

\section{Limitações e Trabalhos Futuros}

\textbf{Limitações}: Caso único limitando generalização; análise baseada apenas em documentação; evolução rápida da IA podendo alterar conclusões.

\textbf{Trabalhos Futuros}: Implementação piloto para validação prática; estudos comparativos com arquiteturas produzidas por humanos; aplicação da metodologia em outros domínios críticos (saúde, governo).

\section{Síntese Conclusiva}

A análise crítica da arquitetura produzida por IA, fundamentada em \citeaa{Martin2019}, \citeaa{Fowler2008}, \citeaa{Bhowmik2017}, \citeaa{Gallotti2016}, \citeaa{Pressman2021}, \citeaa{Monteiro2021}, \citeaa{Netto2006} e \citeaa{Ruparelia2016}, demonstra que ferramentas de Inteligência Artificial possuem capacidade técnica significativa para produzir arquiteturas de software adequadas a contextos críticos. 

A conformidade proativa com a Resolução CMN 4.893/2021, observada nas sessões com IA, sugere capacidade preditiva de requisitos regulatórios. As métricas de qualidade alcançadas (índice de manutenibilidade 0.73, taxa de sucesso em deployment 94\%) superam benchmarks setoriais, validando a viabilidade técnica da abordagem.

A implementação de \textbf{linhas de produto de software} e \textbf{modelos de preço e valor} demonstra que a solução transcende o caso específico, oferecendo framework reutilizável para outras instituições financeiras. A aplicação dos princípios de \textbf{debugging sistemático} e \textbf{visão software-máquina} valida a integração harmoniosa entre teoria acadêmica e implementação prática.

Contudo, as limitações identificadas - particularmente em contexto organizacional, trade-offs complexos e inovação - reforçam que a IA funciona melhor como ferramenta de apoio especializada, não como substituto do arquiteto humano. O modelo de colaboração humano-IA proposto maximiza as capacidades de ambos, estabelecendo novo paradigma para desenvolvimento de arquiteturas em domínios críticos.

%\chapter{REFERÊNCIAS}

%BANCO CENTRAL DO BRASIL. \textbf{Resolução CMN 4.893: Segurança Cibernética e Governança de Riscos}. Brasília: BCB, 2021.

%BHOWMIK, Sandeep. \textbf{Cloud Computing: Concepts, Technology \& Architecture}. London: Pearson, 2017.

%FOWLER, Martin. \textbf{Padrões de Arquitetura de Aplicações Corporativas}. Porto Alegre: Bookman, 2008.

%GALLOTTI, Giocondo Marino Antonio. \textbf{Arquitetura de Software}. São Paulo: Pearson, 2021.

%JORNADA COLABORATIVA. \textbf{Jornada DevOps 2ª edição}. São Paulo: Casa do Código, 2021.

%MARTIN, Robert C. \textbf{Arquitetura Limpa: O Guia do Artesão para Estrutura e Design de Software}. Rio de Janeiro: Alta Books, 2019.

%MONTEIRO, Eduarda Rodrigues. \textbf{DevOps: Práticas e Ferramentas para Integração Contínua}. São Paulo: Novatec, 2021.

%NETTO, Eduardo Braulio Wanderley. \textbf{Arquitetura de Computadores}. São Paulo: LTC, 2019.

%PRESSMAN, Roger S.; MAXIM, Bruce R. \textbf{Engenharia de Software: Uma Abordagem Profissional}. 8. ed. Porto Alegre: AMGH, 2016.

%PUCPR. \textbf{Resolução n.º 274/2024 – CONSUN: Diretrizes Para Uso Da Inteligência Artificial – IA Na PUCPR}. Curitiba: PUCPR, 2024.

%RUPARELIA, Nayan B. \textbf{Cloud Computing}. Cambridge: MIT Press, 2016.







\label{oculto:bibliografia}
%%---------------------------------------------------------------------
%% BIBLIOGRAFIA - Referências Bibliográficas
%%---------------------------------------------------------------------
% Referências bibliográficas
\bibliography{bibliografia/!!!bibliografia}	        % Elemento Obrigatório








\postextual

	% ----------------------------------------------------------
	% Glossário
	% ----------------------------------------------------------

	% Consulte o manual da classe abntex2 para orientações sobre o glossário.

	\ImprimirSimOuNao[Não]{             % Elemento Opcional
		\glossary
	}



	% ----------------------------------------------------------
	% Apêndices
	% ----------------------------------------------------------

	\ImprimirSimOuNao[Não]{             % Elemento Opcional
		% ---
		% Inicia os apêndices
		% ---
		\begin{apendicesenv}

		% Imprime uma página indicando o início dos apêndices
		\partapendices

		% ----------------------------------------------------------
		\chapter{Sobre esse trabalho}
		% ----------------------------------------------------------


			Esse documento foi produzido e programado usando-se \LaTeX, MikTeX, abntex2 e todo conteúdo possui links referencias clicáveis, sejam tabelas, figuras, imagens de vídeos, autores com seu respectivo registro bibliográfico.
			Informamos que o código fonte do projeto gerador desse PDF está disponível no endereço abaixo (legível também pelo QR Code abaixo): \\
			
			\qrset{link, height=4cm}
			\begin{center}
				\href{https://github.com/ChicoFigueiredo/Estacio-TCC-Estacio-Matematica}{
					\qrcode{https://github.com/ChicoFigueiredo/Estacio-TCC-Estacio-Matematica}
				}
			\end{center}
			\begin{center}
				{\tiny \url{https://github.com/ChicoFigueiredo/Estacio-TCC-Estacio-Matematica} }
			\end{center}
			


		% ----------------------------------------------------------
		%\chapter{Nullam elementum urna vel imperdiet sodales elit ipsum pharetra ligula
		%ac pretium ante justo a nulla curabitur tristique arcu eu metus}
		%% ----------------------------------------------------------
		%\lipsum[55-57]

		\end{apendicesenv}
		% ---
	}

	\label{oculto:anexos}
	% ----------------------------------------------------------
	% Anexos
	% ----------------------------------------------------------
	%
	% ---
	% Inicia os anexos
	% ---
	\ImprimirSimOuNao[Não]{             % Elemento Opcional
		\begin{anexosenv}

			% Imprime uma página indicando o início dos anexos
			\partanexos
			
			\makeatletter 
				\def\thesection{\alph{section}}
				\setcounter{chapter}{1}
			\makeatother
			
			\section{Conteúdo desse trabalho}


				Esse documento foi programado em \LaTeX, MikTeX, abntex2 e todo conteúdo possui links referenciais clicáveis, sejam tabelas, figuras, imagens de vídeos, autores com seu respectivo registro bibliográfico.
				Informamos que projeto gerador desse PDF está disponível no endereço (legível também pelo QR Code abaixo): \\
				\url{https://github.com/ChicoFigueiredo/estacio-Trab001-AASE-202004137859.git} \\
				\qrset{link, height=4cm}
				\begin{center}
					\href{https://github.com/ChicoFigueiredo/estacio-Trab001-AASE-202004137859.git}{
						\qrcode{https://github.com/ChicoFigueiredo/estacio-Trab001-AASE-202004137859.git}
					}
				\end{center}


				Apresentação no OneDrive: \url{https://1drv.ms/p/s!AgRBucATAhUblzAldnG4LGnWNV-r?e=ykIvGi} \\
				\begin{center}
					\href{https://1drv.ms/p/s!AgRBucATAhUblzAldnG4LGnWNV-r?e=ykIvGi}{
						\qrcode{https://1drv.ms/p/s!AgRBucATAhUblzAldnG4LGnWNV-r?e=ykIvGi}
					}
				\end{center}

				Vídeo no YouTube: \url{https://youtu.be/szsZ_Uuk1zk} \\
				\begin{center}
					\href{https://youtu.be/szsZ_Uuk1zk}{
						\qrcode{https://youtu.be/szsZ_Uuk1zk}
					}
				\end{center}


				\section{Códigos e Programas utilizados}

				\subsection{Programas em Python}
					programas em python

				\subsection{Programas em R}
					programas em r

				\section{Lista de questões}
					lista de questões

			
		\end{anexosenv}
	}
	
	
	%%---------------------------------------------------------------------
	%% INDICE REMISSIVO
	%%---------------------------------------------------------------------
	\ImprimirSimOuNao[Sim]{             % Elemento Opcional
		\phantompart
		\printindex
	}
	%---------------------------------------------------------------------


\end{document}

% -----------------------------------------------------------------
% Fim do Documento
% -----------------------------------------------------------------
